% !TeX root = ../../Skript.tex
\cohead{\Large\textbf{Koordinatendarstellung}}
\fakesubsection{Koordinatendarstellung}
Wir verwenden nach wie vor das Kartesische Koordinatensystem, nur nennen wir nun die Achsen \(x_1\)-Achse und \(x_2\)-Achse. Hinweis: Wir könnten auch bei der alten \(x\)-\(y\)-Benennung bleiben, wählen jedoch die neue Benennung, da diese auch so im Abitur Verwendung findet.

Einen Vektor können wir dann als \(\vv{v}=\spvek{v_1}{v_2}\) schreiben, wobei \(v_1\) die Länge des Vektors in \(x_1\)-Richtung und \(v_2\) die Länge des Vektors in \(x_2\)-Richtung ist:

\bigskip

\begin{minipage}{\textwidth}
	\adjustbox{valign=t, margin=0ex 0ex 5ex 0ex}{\begin{minipage}{0.5\textwidth-5ex}\raggedright
			Eingezeichnet ist der Vektor \(\vv{a}\) vom Punkt \((1\vert1)\) zum Punkt \((4\vert2)\). Die Länge des Vektors in \(x_1\)-Richtung beträgt \(a_1=3\) und die Länge in \(x_2\)-Richtung beträgt \(a_2=1\). In der Koordinatendarstellung ergibt sich \[\vv{a}=\spvek{3}{1}\]
			Zeichne die Vektoren \(\vec{b}=\spvek{-1}{4}\), \(\vec{c}=\spvek{-2}{-3}\), \(\vec{d}=\spvek{0}{1}\) und \(\vec{e}=\spvek{2}{-6}\) jeweils von einem beliebigen Punkt aus ein.
	\end{minipage}}%
	\adjustbox{valign=t}{\begin{minipage}{0.5\textwidth}
		\includegraphics[width=\textwidth]{\vektor/pics/Koordinatendarstellung_empty.png}%TODO Bild aufdoppeln
	\end{minipage}}%
\end{minipage}

\begin{Exercise}[title={\raggedright Aufgaben im Schulbuch Seiten 194 und 194}, label=VektorenKoordinatendarstellungA1]
	Von besonderem Interesse sind Aufgaben 1, 2, 3, 4, 7, 8 und 12.
\end{Exercise}
\newpage
%%%%%%%%%%%%%%%%%%%%%%%%%%%%%%%%%%%%%%%%%
\begin{Answer}[ref=VektorenKoordinatendarstellungA1]
 xxx%TODO
\end{Answer}